\clearpage
\fancyhead[L]{\small PALLAS.}
\hypertarget{native-465-pallas}{}
\begin{center}{\large PALLAS.}\end{center}

\begin{center}
\rule{0.55\linewidth}{0.4pt}\\[2pt]
Astronomische Nachrichten. Vol. VI. No. 124. cols. 67\,..\,68. October 1827.\\[2pt]
\rule{0.55\linewidth}{0.4pt}
\end{center}

\begin{center}Observations of Pallas.\end{center}

Councillor \textsc{Gauss} shared with me the following observations he made of Pallas.

\begin{center}
\begin{tabular}{l|l|l|l}
 & & R.A. & $\delta$ \\
\hline
1827 Aug. 31 & $11^{\rm u}16^{\rm m}2^{\rm s}0$ & $328^\circ19'54''6$ & $+5^\circ52'27''4$ \\
\hfill Sept. 2 & $11\;6\;42.2$ & $327\;57\;51.4$ & $5\;28\;38.2$ \\
\hfill 3 & $11\;2\;2.8$ & $327\;46\;57.4$ & $5\;16\;31.8$ \\
\hfill 4 & $10\;57\;24.2$ & $327\;36\;14.9$ & $5\;4\;21.3$ \\
\end{tabular}
\end{center}

The first observation is less reliable because the large difference between the temperatures outside and inside made the planet extremely difficult to see. For the other 3, the windows had been opened long beforehand, and Councillor \textsc{Gauss} considers these observations very good.
\begin{flushright}\textsc{Schumacher.}\end{flushright}

